\documentclass[12pt]{article}
\usepackage[utf8]{inputenc}
\usepackage[T1]{fontenc}
\usepackage{hyperref}
\title{Harappan Language Resources for K-12 Curriculum Development}
\author{Compiled by Claude Sonnet 5.5 (claude-sonnet-5-5) for Andr\'as Kornai}
\date{}
\begin{document}
\maketitle

\section{Introduction}
Harappan (ISO 639-3 xiv; Glottocode hara1272) is the presumed language of the Indus Valley (Harappan) civilisation. It is unattested in any decipherable form. The Indus script is undeciphered, and it has not been demonstrated to be a writing system, so the language is unknown. Proposals include proto-Dravidian (Asko Parpola is a notable proponent), Sanskrit, and Munda languages such as Ho and Santali, often on scant evidence. Inscriptions are very short, averaging about 5 signs with a maximum of 26. No bilingual text exists. Sign-count estimates range from 62 to 676.

\section{Primary Educational Resources}
There are no native speakers and no language to teach. The search found no primers or textbooks. A curriculum that uses the language for as much as possible is not possible. At most, schools could teach about the Indus script and the undeciphered-writing problem, in history, archaeology or mathematics classes.

\section{Linguistic and Reference Materials}
\begin{itemize}
\item Wikipedia, ``Harappan language'': \url{https://en.wikipedia.org/wiki/Harappan_language}
\item Harappa.com, ``Why has the language of Harappa not been deciphered'': \url{https://www.harappa.com/answers/why-has-language-harappa-not-been-deciphered}
\item Compass (RAU's IAS), ``Why have scholars struggled to decipher Harappan Script?'': \url{https://compass.rauias.com/current-affairs/scholars-struggled-decipher-harappan-script/}
\item IAS Gyan, ``Deciphering Harappan Script'': \url{https://www.iasgyan.in/daily-current-affairs/deciphering-harappan-script}
\end{itemize}

\section{Religious and Community Resources}
None. There is no known community of speakers.

\section{Digital and Audio Resources}
There is no audio. The search found no corpus resources, apart from mentions of AI decipherment attempts and a decipherment prize.

\section{Curriculum Development Context}
A heritage or revitalization curriculum is not feasible for an unattested language. Any school use would be cultural and historical, presenting the civilisation and the open decipherment problem honestly. It should not present any proposed decipherment as established.
\end{document}